\par\Needspace{7\baselineskip}\bigskip
\begin{center}
\pdfbookmark[2]{Gauss to Olbers, Göttingen, 31 December 1814}{pallas-brief-18141231}
\textsc{Gauss} to \textsc{Olbers}, Göttingen, 31 December 1814.
\end{center}
\noindent\hspace*{\parindent}. . . . .\quad You have probably read a brief notice about the latest opposition of Pallas in our Gel. Anz. Since then I have fitted the elements more closely still and find the following differences:
\begin{center}
\begingroup\small\setlength{\tabcolsep}{5pt}
\begin{tabular}{@{}lrrrr@{}}
&\multicolumn{2}{c}{Old mass of Jupiter \(\left[\frac{1}{1067{,}09}\right]\)}&\multicolumn{2}{c}{Mass of Jupiter increased by \(\frac{1}{64}\)}\\[3pt]
&\multicolumn{1}{c}{Mean long.}&\multicolumn{1}{c}{Hel. lat.}&\multicolumn{1}{c}{Mean long.}&\multicolumn{1}{c}{Hel. lat.}\\
I.&\(-2''{,}6\)&\(-5''{,}1\)&\(+4''{,}4\)&\(-3''{,}0\)\\
II.&\(-11{,}2\)&\(-12{,}2\)&\(+1{,}8\)&\(-8{,}5\)\\
III.&\(-17{,}3\)&\(-0{,}5\)&\(-5{,}9\)&\(-3{,}3\)\\
IV.&\(+3{,}2\)&\(-6{,}7\)&\(-10{,}8\)&\(-6{,}1\)\\
V.&\(+22{,}8\)&\(+9{,}0\)&\(+3{,}3\)&\(+11{,}5\)\\
VI.&\(+23{,}8\)&\(-3{,}4\)&\(+1{,}1\)&\(-0{,}5\)\\
VII.&\(+16{,}8\)&\(-7{,}2\)&\(+13{,}2\)&\(-5{,}4\)\\
VIII.&\(-17{,}9\)&\(+4{,}5\)&\(-1{,}2\)&\(+2{,}2\)\\
IX.&\(-9{,}6\)&\(-0{,}1\)&\(-4{,}0\)&\(-0{,}7\)\\
X.&\(-7{,}7\)&\(+2{,}2\)&\(-2{,}0\)&\(+5{,}5\)\\
Sum of sq.&2284&391&381&325.
\end{tabular}
\endgroup
\end{center}
\par\SourceLeaf{432}Here, however, only the perturbations by Jupiter as calculated by what I call the \textls[180]{special} calculation (carried forward from one year to the next by mechanical quadrature) [are taken as the basis]. The calculation by my general formulas, as well as that of the perturbations by Saturn, is still unfinished. It takes several months each time. I shall soon work out my method and make it known in one way or another, so as finally to be rid of a task whose good results have rewarded me, but which in itself is all too burdensome.
\par
